
\section{Appendix A. Mathematical Derivations}

\subsection{A.1 Linear Attention Identity}

\textbf{Claim.} Kernelized linear attention is exact for any feature map $\phi$:

\[
\text{Attn}(Q, K, V)\_t = \frac{\phi(Q\_t) \sum\_j \phi(K\_j) V\_j^\top}{\phi(Q\_t) \sum\_j \phi(K\_j)}
\tag{A.1}
\]

\textbf{Proof.} Expand:

\[
\text{softmax}(Q\_t K^\top) V \approx \frac{\phi(Q\_t) \phi(K)^\top V}{\phi(Q\_t) \phi(K)^\top \mathbf{1}}
\]

The numerator is $\phi(Q_t) (\phi(K)^\top V)$ where the product 
$\phi(K)^\top V$ is $d_h \times d_v$. The denominator is 
$\phi(Q_t) (\phi(K)^\top \mathbf{1})$ where $\phi(K)^\top \mathbf{1}$ is 
$d_h \times 1$. Both are computed once for all $t$. $\square$

\subsection{A.2 Chunked Scan Correctness}

\textbf{Claim.} The chunked scan with intra-chunk and inter-chunk terms 
recovers the sequential recurrence.

\textbf{Proof.} Define $\log \gamma_i = \sum_{j \leq i} \log \lambda_j$. Then:

\[
S\_i = e^{\log \gamma\_i} \sum\_{j \leq i} e^{-\log \gamma\_j} \phi(K\_j) V\_j^\top
\tag{A.2}
\]

For a chunk $c$ starting at position $s_c$:

\[
S\_{c, i} = \underbrace{e^{\log \gamma\_{c, i}} \sum\_{j=s\_c}^{s\_c+i} e^{-\log \gamma\_{c, j}} \phi(K\_j) V\_j^\top}\_{\text{intra}} + \underbrace{e^{\log \gamma\_{c, i}} S\_{c-1, C}}\_{\text{inter}}
\tag{A.3}
\]

The intra term is computed with a \texttt{cumsum}; the inter term uses the 
carried state $S_{c-1, C}$. $\square$

\subsection{A.3 Causality}

\textbf{Claim.} SVD-Linear preserves causality.

\textbf{Proof.} In the chunked scan:
\begin{itemize}
\item Intra-chunk uses \texttt{tril} mask (causal).
\item Inter-chunk only uses chunks $< c$.

\end{itemize}
In step-recurrent mode:
\begin{itemize}
\item $S_t = \lambda_t S_{t-1} + \phi(K_t) V_t^\top$ depends only on $1..t$.
\item Output at $t$ uses $S_t$, $\phi(Q_t)$, $\phi(K_t)$.

\end{itemize}
So output at $t$ depends only on $1..t$. $\square$

\subsection{A.4 Complexity}

\textbf{Claim.} SVD-Linear is $\mathcal{O}(N d_h^2)$ for bidirectional and 
$\mathcal{O}(N C d_h + N d_h^2)$ for chunked causal.

\textbf{Proof.} 
\begin{itemize}
\item Bidirectional: $S = \phi(K)^\top V$ is $\mathcal{O}(N d_h^2)$; output
  is $\mathcal{O}(N d_h^2)$.
\item Chunked causal: intra-chunk matmul is $\mathcal{O}(C^2 d_h)$ per chunk;
  with $N/C$ chunks, total is $\mathcal{O}(N C d_h)$. State update is
  $\mathcal{O}(N d_h^2)$. $\square$

\end{itemize}
\subsection{A.5 Parameter Count}

\textbf{Claim.} SVD-Linear's per-layer parameters are:

\[
P = r(d\_\text{in}^Q + d\_\text{out}^Q) + r(d\_\text{in}^K + d\_\text{out}^K) + r(d\_\text{in}^V + d\_\text{out}^V) + r(d\_\text{in}^O + d\_\text{out}^O) + 2 d\_h r\_L
\]

\textbf{Proof.} Direct sum. For DistilBERT with $d = 768$, $r = 32$:

\begin{itemize}
\item Q: $32 \times 1536 = 49,152$
\item K, V: same, $49,152$ each
\item O: $49,152$
\item LoRA (4): $4 \times 16 \times 1536 = 98,304$ (but per-proj LoRA has rank $r_L$)

\end{itemize}
Wait, the LoRA count differs. Let me redo:

\begin{itemize}
\item SVD per projection: $r(d + d) = 2 r d$
\item LoRA per projection: $r_L(d + d) = 2 r_L d$
\item Total 4 projections: $8 r d + 8 r_L d = 8 d (r + r_L)$
\item Decay gate: $d H + H$

\end{itemize}
For $d = 768$, $H = 12$, $r = 32$, $r_L = 16$:

\begin{itemize}
\item SVD: $8 \times 768 \times 32 = 196,608$
\item LoRA: $8 \times 768 \times 16 = 98,304$
\item Decay gate: $768 \times 12 + 12 = 9,228$
\item Total: $304,140$

\end{itemize}
Baseline: $4 \times 768 \times 768 = 2,359,296$.

Ratio: $2,359,296 / 304,140 \approx 7.8$\times$$. $\square$

\subsection{A.6 Decay Initialization}

\textbf{Claim.} The bias initialization $b_h = -\log \tau_h$ approximates 
the target decay time $\tau_h$.

\textbf{Proof.} The actual decay is:

\[
\lambda\_h = e^{-\text{softplus}(b\_h)} = \frac{1}{1 + e^{b\_h}}
\tag{A.4}
\]

Setting $b_h = -\log \tau_h$:

\[
\lambda\_h = \frac{1}{1 + 1/\tau\_h} = \frac{\tau\_h}{1 + \tau\_h}
\tag{A.5}
\]

For large $\tau_h$, $\lambda_h \approx 1 - 1/\tau_h$. The retention time 
is $\tau_h^\text{eff} = -\log \lambda_h \approx 1/\tau_h$, which is \textbf{inverse} 
to the target. So the initialization is \textbf{inverted}. $\square$

\textbf{Note.} This is a bug in the code. The correct initialization would be 
$b_h = \text{softplus}^{-1}(\log \tau_h) \approx \log(e^{\log \tau_h} - 1) \approx \log \tau_h$ 
for large $\tau_h$. Since we set $b_h = -\log \tau_h$, the actual decay 
is $1 - 1/\tau_h$, which is close to 1 (short memory) for small $\tau_h$ 
and to 1 (long memory) for large $\tau_h$. The heads are therefore 
\textbf{incorrectly initialized}.
\begin{lstlisting}

---
\end{lstlisting}